Data source and setting: single-centre cohort from ASUIT. {[}Type of data: CTA, clinical follow-up, CFD.]

\section{Cohort and outcome definition}
Items 6--8 of TRIPOD+AI \cite{collins_tripodai_2024}.
\begin{itemize}
    \item Inclusion and exclusion criteria. {[}To define.]
    \item Minimum follow-up: the only outcome-related exclusion; it avoids false negatives.
    \item Flowchart of included and excluded patients.
    \item Primary outcome $D \in \{0,1\}$ and time horizon. {[}To define with clinicians.]
\end{itemize}

\section{Data preparation: from CTA to CFD simulation}
\begin{enumerate}
    \item From CTA to CFD simulation pipeline.  \cite{lasso_slicerheart_2022, trentin_automatic_2015}.
    \item Solver setup (svMultiPhysics) and boundary conditions: inflow and RCR outlets \cite{du_outflow_2015, ramazanli_modeling_2025, les_quantification_2010}. \citeneeded{svMultiPhysics}
    \item Quality control of the simulations and excluded patients (Appendix~\ref{app:cfd}).
\end{enumerate}
Modelling assumptions: rigid wall, Newtonian fluid, modelled outlets.

\section{Features}

\subsection{Morphological and radiomic features}
\begin{itemize}
    \item Morphological and radiomic features, as defined by the reference model (Appendix~\ref{app:benchmark}).
\end{itemize}

\subsection{Hemodynamic features}
\begin{itemize}
    \item Indices from Section~\ref{sec:wss-indices}, on regions $R_k$: proximal neck, sac, iliacs. regions are the same as the ones used for morphology extraction in the benchmark. 
\end{itemize}
\begin{equation}
    \bar{q}_{k} = \frac{1}{|R_k|}\int_{R_k} q(\mathbf{x})\,dA,
    \qquad
    f_{k}(\theta) = \frac{1}{|R_k|}\int_{R_k} \mathbb{1}\left[q(\mathbf{x}) > \theta\right] dA
\end{equation}
\begin{itemize}
    \item Pre-specified set $\mathbf{Hd}_i \in \mathbb{R}^p$, with small $p$. 
    \item idea: use as feature the portion of the surface of the region that surpasses a thresholds $\theta$ identified from the literature (e.g.: LOW TAWSS<0.4 Pa)

    \item in the more complex representation (pointnet++), we keep a common anatomical domain for all patients (from proximal neck to common iliacs). we keep the representation of the hemodynamic field as a point-like representation
    \item $N = 4096$ wall points per patient, sampled with probability proportional to area; resampled at each training epoch.
    \item Per point: coordinates and $\mathbf{q} = (\log\mathrm{TAWSS}, \mathrm{OSI}, \log\mathrm{RRT})$, standardised.
\end{itemize}
\begin{equation}
    \mathcal{P}_i = \left\{ \left(\mathbf{x}_{ij},\, \mathbf{q}_{ij}\right) \right\}_{j=1}^{N}
\end{equation}

